\documentclass[10pt,a4paper]{amsart}
\usepackage{amsmath,amssymb,amsthm,graphicx}
\newcommand{\RS}[1]{\textcolor{red}{#1}}

%\usepackage{refcheck}

\numberwithin{equation}{section}

\newtheorem{Theorem}{Theorem}[section]
\newtheorem*{Theorem*}{Theorem}
\newtheorem{Corollary}[Theorem]{Corollary}
\newtheorem{Lemma}[Theorem]{Lemma}
\newtheorem{Proposition}[Theorem]{Proposition}
 { \theoremstyle{definition}
\newtheorem{Definition}[Theorem]{Definition}
\newtheorem{Note}[Theorem]{Note}
\newtheorem{Example}[Theorem]{Example}
\newtheorem{Remark}[Theorem]{Remark} }

\newcommand{\sech}{\operatorname{sech}}
\usepackage{cleveref}
\crefname{originalequation}{original equation}{original equations}
\Crefname{originalequation}{Original equation}{Original equations}
\makeatletter
\newlabel{eq:rs}{{internal}{source}{}{sourcepointer-1}{}}
\newlabel{eq:rs@cref}{{[originalequation][0][]internal}{[1][0][]source}{}{}{}}
\makeatother
\crefname{originalequation}{original equation}{original equations}
\Crefname{originalequation}{Original equation}{Original equations}
\makeatletter
\newlabel{eq:epphi}{{internal}{source}{}{sourcepointer-2}{}}
\newlabel{eq:epphi@cref}{{[originalequation][0][]internal}{[1][0][]source}{}{}{}}
\makeatother
\crefname{originaltheorem}{original theorem}{original theorems}
\Crefname{originaltheorem}{Original theorem}{Original theorems}
\makeatletter
\newlabel{thBT}{{internal}{source}{}{sourcepointer-3}{}}
\newlabel{thBT@cref}{{[originaltheorem][0][]internal}{[1][0][]source}{}{}{}}
\makeatother
\crefname{originalequation}{original equation}{original equations}
\Crefname{originalequation}{Original equation}{Original equations}
\makeatletter
\newlabel{eqorig1_4}{{internal}{source}{}{sourcepointer-4}{}}
\newlabel{eqorig1_4@cref}{{[originalequation][0][]internal}{[1][0][]source}{}{}{}}
\makeatother
\crefname{originalequation}{original equation}{original equations}
\Crefname{originalequation}{Original equation}{Original equations}
\makeatletter
\newlabel{eq:btanalytic_vc1_4}{{internal}{source}{}{sourcepointer-5}{}}
\newlabel{eq:btanalytic_vc1_4@cref}{{[originalequation][0][]internal}{[1][0][]source}{}{}{}}
\makeatother
\title{Superposition formula for two Backlund transformations of the first family}
\author{Kelmer, Filipe, Tenenblat, Keti}
\date{doi:10.3842/SIGMA.2024.015; CC BY 4.0}
\begin{document}
\maketitle
\setcounter{section}{3}
\setcounter{Theorem}{0}
\noindent\textit{Editorial note.} All delivered mathematical text is the author's own, copied from the pinned SIGMA article source; no statement or proof step was rewritten, reordered or removed. The counted statement is the labeled theorem environment at source lines 362-376 and the counted proof is the proof environment at source lines 379-626.\par\medskip
\noindent\textit{Editorial note.} The retained preamble of this unit was generated mechanically for this delivery: it copies the author's own numbering, macro and theorem declaration block from the pinned source (source lines 9-27), with the publisher class and page setup lines removed.\par\medskip
\noindent\textit{Allowed standard prerequisites (editorial).}\begin{itemize}
\item The analytic Backlund transformation relating a solution to its companion at a parameter. same article, the corollary at source lines 308-317. Fixes the equations satisfied by the two companion solutions and the parameters, which are the input of the algebraic computation.
\item The Backlund transformation theorem of the first family with its two branches. same article, the theorem at source lines 270-285. Supplies the general form of the transformation, from which the exhibited superposition formula is derived.
\item The two admissible parameter branches encoded by the pair of integrality parameters. same article, the displays at source lines 191 and 199. Restricts the parameters so that the transformations are of the first family, which is the hypothesis of the statement.
\end{itemize}\medskip
\ifdefined\hypertarget\hypertarget{sourcepointer-1}{}\fi
\noindent\textit{Source reference.} Original equation internal: sources/sigma-2024-015-superposition-formulae-for-the-geometric-b-2/pinned-original-source.tex, source lines 191--191. The counted statement or proof refers to that displayed integrality condition on the parameters; it is the author's own earlier display or statement in the same article and the whole argument of the counted span is retained verbatim.\par\medskip
\ifdefined\hypertarget\hypertarget{sourcepointer-2}{}\fi
\noindent\textit{Source reference.} Original equation internal: sources/sigma-2024-015-superposition-formulae-for-the-geometric-b-2/pinned-original-source.tex, source lines 199--199. The counted statement or proof refers to that displayed relation between the two transformation parameters; it is the author's own earlier display or statement in the same article and the whole argument of the counted span is retained verbatim.\par\medskip
\ifdefined\hypertarget\hypertarget{sourcepointer-3}{}\fi
\noindent\textit{Source reference.} Original theorem internal: sources/sigma-2024-015-superposition-formulae-for-the-geometric-b-2/pinned-original-source.tex, source lines 270--285. The counted statement or proof refers to that Backlund transformation theorem of the first family; it is the author's own earlier display or statement in the same article and the whole argument of the counted span is retained verbatim.\par\medskip
\ifdefined\hypertarget\hypertarget{sourcepointer-4}{}\fi
\noindent\textit{Source reference.} Original equation internal: sources/sigma-2024-015-superposition-formulae-for-the-geometric-b-2/pinned-original-source.tex, source lines 308--313. The counted statement or proof refers to that displayed equation satisfied by the companion solutions; it is the author's own earlier display or statement in the same article and the whole argument of the counted span is retained verbatim.\par\medskip
\ifdefined\hypertarget\hypertarget{sourcepointer-5}{}\fi
\noindent\textit{Source reference.} Original equation internal: sources/sigma-2024-015-superposition-formulae-for-the-geometric-b-2/pinned-original-source.tex, source lines 314--317. The counted statement or proof refers to that displayed analytic Backlund transformation of the first family; it is the author's own earlier display or statement in the same article and the whole argument of the counted span is retained verbatim.\par\medskip
\begin{Theorem}\label{th:superth1_4}
 Consider fixed constants $\delta=\pm 1$, $\epsilon=\pm 1$,
$r,s\in \{0,1\}$, $r\leq s$, satisfying~\eqref{eq:rs} and \eqref{eq:epphi}, and $\delta=(-1)^{s+1}$. Let $\alpha(x_1,x_2)$ be a solution of \eqref{eqorig1_4} and
 let $\alpha'$ $($resp.\ $\alpha'')$ be a~solution of \eqref{eqorig1_4}
 associated to $\alpha$ by the $\phi_1($resp.\ $\phi_2)$-B\"acklund transformation \eqref{eq:btanalytic_vc1_4}.
 If $\phi_1\neq \phi_2$, then there exists a unique solution $\alpha^*$ of
 \eqref{eqorig1_4} such that $\alpha^*$ is associated to $\alpha'$ $($resp.~$\alpha'')$ by the $\phi_2($resp.\ $\phi_1)$-B\"acklund transformation \eqref{eq:btanalytic_vc1_4}. Moreover, $\alpha^*$ is
given algebraically by
\begin{equation}\label{eq:superpos}
T_{(-1)^r}\biggl(\frac{\alpha^*-\alpha}{4}\biggr) =
\delta\tau \frac{S_{-\delta\epsilon}\bigl(\frac{\phi_2+\phi_1}{2}\bigr) } {S_{-\delta\epsilon}\bigl(\frac{\phi_2-\phi_1}{2}\bigr)}
T_{(-1)^r}\biggl(\frac{\alpha'-\alpha''}{4}\biggr),
\end{equation}
where $T_{(-1)^r}$ is the function $\tan$ $($resp.\ $\tanh)$ when $r=0$ $($resp.\ $r=1)$.
\end{Theorem}

\begin{proof}
Let $\alpha(x_1,x_2)$ be a solution of \eqref{eqorig1_4}. Let $\alpha'$ (resp.\ $\alpha''$) be a solution of \eqref{eqorig1_4} associated to $\alpha$ by the B\"acklund transformation \eqref{eq:btanalytic_vc1_4}
for parameters $\{\lambda_1=S_{-\epsilon\delta}(\phi_1),\; \Lambda_1=C_{-\epsilon\delta}(\phi_1)\}$ (resp.\ $\{\lambda_2=S_{-\epsilon\delta}(\phi_2),\; \Lambda_2=C_{-\epsilon\delta}(\phi_2)\}$), i.e., $\alpha'$ and $\alpha''$ satisfy
\begin{gather}
\alpha'_{x_1}+\alpha_{x_2}=\frac{2}{\lambda_1}S_{l}\biggl(\frac{\alpha'}{2}\biggr)C_{l}\Bigl(\frac{\alpha}{2}\Bigr)-\frac{2(-1)^s\tau\Lambda_1}{\lambda_1} C_{l}\biggl(\frac{\alpha'}{2}\biggr)S_{l}\Bigl(\frac{\alpha}{2}\Bigr),\nonumber\\[1mm]
\alpha'_{x_2}+\alpha_{x_1}=-\frac{2}{\lambda_1}C_{l}\biggl(\frac{\alpha'}{2}\biggr)S_{l}\Bigl(\frac{\alpha}{2}\Bigr)-\frac{2\delta\tau\Lambda_1}{\lambda_1} S_{l}\biggl(\frac{\alpha'}{2}\biggr)C_{l}\Bigl(\frac{\alpha}{2}\Bigr),\label{eq:btanalytic_v1}
\\[1mm]
\alpha''_{x_1}+\alpha_{x_2}=\frac{2}{\lambda_2}S_{l}\biggl(\frac{\alpha''}{2}\biggr)C_{l}\Bigl(\frac{\alpha}{2}\Bigr) -\frac{2(-1)^s\tau\Lambda_2}{\lambda_2}C_{l}\biggl(\frac{\alpha''}{2}\biggr)S_{l}\Bigl(\frac{\alpha}{2}\Bigr),\nonumber\\[1mm]
\alpha''_{x_2}+\alpha_{x_1}=-\frac{2}{\lambda_2}C_{l}\biggl(\frac{\alpha''}
{2 }\biggr)S_{l}\Bigl(\frac{\alpha}{2}\Bigr)-\frac{2\delta\tau\Lambda_2}{\lambda_2} S_{l}\biggl(\frac{\alpha''}{2}\biggr)C_{l}\Bigl(\frac{\alpha}{2}\Bigr),\label{eq:btanalytic_v2}
\end{gather}
 where $ l=(-1)^r$.

Assume there exists $\alpha^*$ associated to $\alpha'$ (resp.\ $\alpha''$) by the B\"acklund transformation with parameters
 $\{\lambda_2,\; \Lambda_2\}$ (resp.\ $\{\lambda_1,\; \Lambda_1\}$). then $\alpha^*$ should satisfy the following systems of equations
\begin{gather}
\alpha^{*}_{x_1}+\alpha'_{x_2} = \frac{2}{\lambda_2}S_{l}\biggl(\frac{\alpha^{*}}{2}\biggr)C_{l}\biggl(\frac{\alpha'}{2}\biggr)-\frac{2(-1)^s\tau\Lambda_2}{\lambda_2} C_{l}\biggl(\frac{\alpha^{*}}{2}\biggr)S_{l}\biggl(\frac{\alpha'}{2}\biggr),\nonumber\\
\alpha^{*}_{x_2}+\alpha'_{x_1}=-\frac{2}{\lambda_2}C_{l}\biggl(\frac{\alpha^{*}}{2}\biggr)S_{l}\biggl(\frac{\alpha'}{2}\biggr)-\frac{2\delta\tau\Lambda_2}{\lambda_2} S_{l}\biggl(\frac{\alpha^{*}}{2}\biggr)C_{l}\biggl(\frac{\alpha'}{2}\biggr),\label{eq:btanalytic_v2*}\\
\alpha^*_{x_1}+\alpha''_{x_2}=\frac{2}{\lambda_1}S_{l}\biggl(\frac{\alpha^*}{2}\biggr)C_{l}\biggl(\frac{\alpha''}{2}\biggr)-\frac{2(-1)^s\tau\Lambda_1}{\lambda_1} C_{l}\biggl(\frac{\alpha^*}{2}\biggr)S_{l}\biggl(\frac{\alpha''}{2}\biggr),\nonumber\\
\alpha^*_{x_2}+\alpha''_{x_1}=-\frac{2}{\lambda_1}C_{l}\biggl(\frac{\alpha^*}{2}\biggr)S_{l}\biggl(\frac{\alpha''}{2}\biggr)-\frac{2\delta\tau\Lambda_1}{\lambda_1} S_{l}\biggl(\frac{\alpha^*}{2}\biggr)C_{l}\biggl(\frac{\alpha''}{2}\biggr).\label{eq:btanalytic_v1*}
\end{gather}

Subtracting the first (resp.\ second) equation of \eqref{eq:btanalytic_v2} from the first
(resp.\ second) equation of~\eqref{eq:btanalytic_v1}, we obtain
\begin{gather}
\alpha'_{x_1}-\alpha''_{x_1}= \frac{2}{\lambda_1}S_{l}\biggl(\frac{\alpha'}{2}\biggr)C_{l}\Bigl(\frac{\alpha}{2}\Bigr)-\frac{2(-1)^s\tau\Lambda_1}{\lambda_1} C_{l}\biggl(\frac{\alpha'}{2}\biggr)S_{l}\Bigl(\frac{\alpha}{2}\Bigr) -\frac{2}{\lambda_2}S_{l}\biggl(\frac{\alpha''}{2}\biggr)C_{l}\Bigl(\frac{\alpha}{2}\Bigr)
\nonumber\\ \hphantom{\alpha'_{x_1}-\alpha''_{x_1}=}{}
+\frac{2(-1)^s\tau\Lambda_2}{\lambda_2}C_{l}
\biggl(\frac{\alpha''}{2}\biggr)S_{l}\Bigl(\frac{\alpha}{2}\Bigr),\label{eq:primv1-v2}
\\
\alpha'_{x_2}-\alpha''_{x_2}=
-\frac{2}{\lambda_1}C_{l}\biggl(\frac{\alpha'}{2}\biggr)S_{l}\Bigl(\frac{\alpha}{2}\Bigr)-\frac{2\delta\tau\Lambda_1}{\lambda_1} S_{l}\biggl(\frac{\alpha'}{2}\biggr)C_{l}\Bigl(\frac{\alpha}{2}\Bigr) +\frac{2}{\lambda_2}C_{l}\biggl(\frac{\alpha''}
{2 }\biggr)S_{l}\Bigl(\frac{\alpha}{2}\Bigr)
\nonumber\\ \hphantom{\alpha'_{x_2}-\alpha''_{x_2}=}{}
+\frac{2\delta\tau\Lambda_2}{\lambda_2} S_{l}\biggl(\frac{\alpha''}{2}\biggr)C_{l}\Bigl(\frac{\alpha}{2}\Bigr).\label{eq:segv1-v2}
\end{gather}

Similarly, subtracting the first (resp.\ second) equation of \eqref{eq:btanalytic_v1*} from the first (resp.\ second) equation of \eqref{eq:btanalytic_v2*}, we obtain
\begin{gather}
\alpha'_{x_2}-\alpha''_{x_2} = \frac{2}{\lambda_2}S_{l}\biggl(\frac{\alpha^*}{2}\biggr)C_{l}\biggl(\frac{\alpha'}{2}\biggr)+ \frac{2\delta\tau\Lambda_2}{\lambda_2} C_{l}\biggl(\frac{\alpha^*}{2}\biggr)S_{l}\biggl(\frac{\alpha'}{2}\biggr)-\frac{2}{\lambda_1}S_{l}\biggl(\frac{\alpha^*}{2}\biggr)C_{l}\biggl(\frac{\alpha''}{2}\biggr)
\nonumber\\ \hphantom{\alpha'_{x_2}-\alpha''_{x_2} =}{}
-\frac{2\delta\tau\Lambda_1}{\lambda_1} C_{l}\biggl(\frac{\alpha^*}{2}\biggr)S_{l}\biggl(\frac{\alpha''}{2}\biggr),\label{eq:primv2*-v1*}
\\
\alpha'_{x_1}-\alpha''_{x_1}=
-\frac{2}{\lambda_2}C_{l}\biggl(\frac{\alpha^*}{2}\biggr)S_{l}\biggl(\frac{\alpha'}{2}\biggr)-\frac{2\delta\tau\Lambda_2}{\lambda_2} S_{l}\biggl(\frac{\alpha^*}{2}\biggr)C_{l}\biggl(\frac{\alpha'}{2}\biggr)+\frac{2}{\lambda_1}C_{l}\biggl(\frac{\alpha^*}{2}\biggr)S_{l}\biggl(\frac{\alpha''}{2}\biggr)
\nonumber\\ \hphantom{\alpha'_{x_1}-\alpha''_{x_1}=}{}
+\frac{2\delta\tau\Lambda_1}{\lambda_1} S_{l}\biggl(\frac{\alpha^*}{2}\biggr)C_{l}\biggl(\frac{\alpha''}{2}\biggr). \label{eq:segv2*-v1*}
\end{gather}
Now comparing equations \eqref{eq:primv1-v2} with \eqref{eq:segv2*-v1*} and multiplying all terms by $\lambda_1\lambda_2$, we get
\begin{gather*} %\label{eq:comp1}
\lambda_1
\biggl[ C_l\biggl(\frac{\alpha^*}{2}\biggr) S_l\biggl(\frac{\alpha'}{2}\biggr)
-S_l\biggl(\frac{\alpha''}{2}\biggr) C_l\Bigl(\frac{\alpha}{2}\Bigr)\biggr]\\
\qquad{}+\lambda_1\Lambda_2\delta\tau
\biggl[ S_l\biggl(\frac{\alpha^*}{2}\biggr) C_l\biggl(\frac{\alpha'}{2}\biggr)
-C_l\biggl(\frac{\alpha''}{2}\biggr) S_l\Bigl(\frac{\alpha}{2}\Bigr)\biggr] \\
\qquad{}+\lambda_2 \biggl[ S_l\biggl(\frac{\alpha'}{2}\biggr) C_l\Bigl(\frac{\alpha}{2}\Bigl)
 - C_l\biggl(\frac{\alpha^*}{2}\biggr) S_l\biggl(\frac{\alpha''}{2}\biggr) \biggr]\\
 \qquad{}+\lambda_2\Lambda_1\delta\tau \biggl[ C_l\biggl(\frac{\alpha'}{2}\biggr) S_l\Bigl(\frac{\alpha}{2}\Bigr)
 - S_l\biggl(\frac{\alpha^*}{2}\biggr) C_l\biggl(\frac{\alpha''}{2}\biggr) \biggr]=0,
\end{gather*}
where we have used the fact that $\delta=(-1)^{s+1}$.

Using trigonometric identities,
%the identity \eqref{id:SlCl} to each term of
the above equation reduces to
\begin{gather*}
\biggl[S_l\biggl(\frac{\alpha^*+\alpha'}{2}\biggr) - S_l\biggl(\frac{\alpha''+\alpha}{2}\biggr) \biggr]
(\lambda_1 +\delta\tau\lambda_1\Lambda_2)\\
 \qquad{} +\biggl[S_l\biggl(\frac{\alpha^*-\alpha'}{2}\biggr) - S_l\biggl(\frac{\alpha+\alpha''}{2}\biggr) \biggr]
(-\lambda_1 +\delta\tau\lambda_1\Lambda_2) \\
 \qquad{} +\biggl[S_l\biggl(\frac{\alpha+\alpha'}{2}\biggr) - S_l\biggl(\frac{\alpha^*+\alpha'}{2}\biggr) \biggr]
(\lambda_2 +\delta\tau\lambda_2\Lambda_1)\\
 \qquad{} +\biggl[S_l\biggl(\frac{\alpha-\alpha'}{2}\biggr) - S_l\biggl(\frac{\alpha^*-\alpha''}{2}\biggr) \biggr]
(-\lambda_2 +\delta\tau\lambda_2\Lambda_1)=0,
\end{gather*}
 which can be rewritten as
 \begin{gather*}
 C_l\biggl(\frac{\alpha^*+\alpha+\alpha'+\alpha''}{4}\biggr) S_l\biggl(\frac{\alpha^*-\alpha+\alpha'-\alpha''}{4}\biggr) (\lambda_1 +\delta\tau\lambda_1\Lambda_2)\\
 \qquad{} + C_l\biggl(\frac{\alpha^*+\alpha-(\alpha'+\alpha'')}{4}\biggr) S_l\biggl(\frac{\alpha^*-\alpha-(\alpha'-\alpha'')}{4}\biggr) (-\lambda_1 +\delta\tau\lambda_1\Lambda_2)\\
 \qquad{} + C_l\biggl(\frac{\alpha^*+\alpha+\alpha'+\alpha''}{4}\biggr) S_l\biggl(\frac{-(\alpha^*-\alpha)+\alpha'-\alpha''}{4}\biggr) (\lambda_2 +\delta\tau\lambda_2\Lambda_1)\\
 \qquad{} + C_l\biggl(\frac{\alpha^*+\alpha-(\alpha'+\alpha'')}{4}\biggr) S_l\biggl(\frac{-(\alpha^*-\alpha)-(\alpha'-\alpha'')}{4}\biggr) (-\lambda_2 +\delta\tau\lambda_2\Lambda_1)=0.
 \end{gather*}

 By applying again trigonometric identities
 %the identities \eqref{id:Samnb} and \eqref{id:Camnb}
 to each term $C_l$ and $S_l$ above, we obtain the equation in terms of $S_l$ and $C_l$ of the expressions
\[
%D*=
\frac{\alpha^*-\alpha}{4}, \qquad
% D'=
\frac{\alpha'-\alpha''}{4},
\qquad
%Sig*=
\frac{\alpha^*+\alpha}{4}, \qquad
%Sig'=
\frac{\alpha'+\alpha''}{4}.
\]
After cancellations it reduces to
\begin{gather*}
\delta\tau C_l\biggl(\frac{\alpha^*+\alpha}{4}\biggr) C_l\biggl(\frac{\alpha'+\alpha''}{4}\biggr)
S_l\biggl(\frac{\alpha^*-\alpha}{4}\biggr)C_l\biggl(\frac{\alpha'-\alpha''}{4}\biggr)(\lambda_1\Lambda_2-\lambda_2\Lambda_1)\\
\qquad{} -l \delta\tau S_l\biggl(\frac{\alpha^*+\alpha}{4}\biggr) S_l\biggl(\frac{\alpha'+\alpha''}{4}\biggr)
C_l\biggl(\frac{\alpha^*-\alpha}{4}\biggr)S_l\biggl(\frac{\alpha'-\alpha''}{4}\biggr)(\lambda_1\Lambda_2+\lambda_2\Lambda_1)\\
\qquad{} + C_l\biggl(\frac{\alpha^*+\alpha}{4}\biggr) C_l\biggl(\frac{\alpha'+\alpha''}{4}\biggr)
C_l\biggl(\frac{\alpha^*-\alpha}{4}\biggr)S_l\biggl(\frac{\alpha'-\alpha''}{4}\biggr) (\lambda_1+\lambda_2)\\
\qquad{} -lS_l\biggl(\frac{\alpha^*+\alpha}{4}\biggr) S_l\biggl(\frac{\alpha'+\alpha''}{4}\biggr)
S_l\biggl(\frac{\alpha^*-\alpha}{4}\biggr)C_l\biggl(\frac{\alpha'-\alpha''}{4}\biggr) (\lambda_1-\lambda_2)=0.
\end{gather*}

Dividing by
\[
C_l\biggl(\frac{\alpha^*+\alpha}{4}\biggr) C_l\biggl(\frac{\alpha'+\alpha''}{4}\biggr)C_l\biggl(\frac{\alpha^*-\alpha}{4}\biggr)C_l\biggl(\frac{\alpha'-\alpha''}{4}\biggr),
\]
and introducing the notation $T_l=\displaystyle \frac{S_l}{C_l}$, the equation above reduces to
\begin{gather}
 \delta\tau T_l\biggl(\frac{\alpha^*-\alpha}{4}\biggr)(\lambda_1\Lambda_2-\lambda_2\Lambda_1)
 -l\delta\tau T_l\biggl(\frac{\alpha^*+\alpha}{4}\biggr)T_l\biggl( \frac{\alpha'+\alpha''}{4}\biggr)
T_l\biggl(\frac{\alpha'-\alpha''}{4}\biggr) (\lambda_1\Lambda_2+\lambda_2\Lambda_1) \nonumber\\
\qquad{} +T_l\biggl(\frac{\alpha'-\alpha''}{4}\biggr)(\lambda_1+\lambda_2) \nonumber\\
\qquad{} - lT_l\biggl(\frac{\alpha^*+\alpha}{4} \biggr)T_l\biggl(\frac{\alpha'+\alpha''}{4}\biggr)T_l\biggl(\frac{\alpha^*-\alpha}{4}\biggr)(\lambda_1-\lambda_2)=0.\label{Tl}
\end{gather}

Since
\[
\Lambda_i=C_{-\delta\epsilon}(\phi_i) , \qquad \lambda_i=S_{-\delta\epsilon}(\phi_i), \qquad i=1,2,
\]
%Moreover, using the identities \eqref{id:SmnS} and \eqref{id:Samnb}
we have that
\begin{gather*}
\lambda_1\Lambda_2\pm\lambda_2\Lambda_1=2
S_{-\delta\epsilon}\biggl(\frac{\phi_1\pm \phi_2}{2}\biggr)C_{-\delta\epsilon}\biggl(\frac{\phi_1\pm \phi_2}{2}\biggr),
\\
\lambda_1\pm\lambda_2= 2 S_{-\delta\epsilon}\biggl(\frac{\phi_1\pm \phi_2}{2}\biggr)
C_{-\delta\epsilon}\biggl(\frac{\phi_1\mp \phi_2}{2}\biggr).
\end{gather*}
Substituting these expressions into \eqref{Tl}, we conclude that the following product vanishes
\begin{gather*}
 \biggl[T_l\biggl(\frac{\alpha^*-\alpha}{4}\biggr) S_{-\delta\epsilon}\biggl(\frac{\phi_2-\phi_1}{2}\biggr)-\delta\tau
T_l\biggl(\frac{\alpha'-\alpha''}{4}\biggr) S_{-\delta\epsilon}\biggl(\frac{\phi_2+\phi_1}{2}\biggr)\biggr] \\
 \qquad{}\times\biggl[ -\delta\tau C_{-\delta\epsilon}\biggl(\frac{\phi_1-\phi_2}{2}\biggr)+l T_l\biggl(\frac{\alpha^*+\alpha}{4}\biggr)
T_l\biggl(\frac{\alpha'+\alpha''}{4}\biggr) C_{-\delta\epsilon}\biggl(\frac{\phi_1+\phi_2}{2}\biggr) \biggr]=0.
\end{gather*}
Similar computations, comparing equations \eqref{eq:segv1-v2} and \eqref{eq:primv2*-v1*}
show that
\begin{gather*}
 \biggl[T_l\biggl(\frac{\alpha^*-\alpha}{4}\biggr) S_{-\delta\epsilon}\biggl(\frac{\phi_2-\phi_1}{2}\biggr)-\delta\tau
T_l\biggl(\frac{\alpha'-\alpha''}{4}\biggr) S_{-\delta\epsilon}\biggl(\frac{\phi_2+\phi_1}{2}\biggr)\biggr] \\
 \qquad{}\times\biggl[ C_{-\delta\epsilon}\biggl(\frac{\phi_1+\phi_2}{2}\biggr)-\delta\tau l T_l\biggl(\frac{\alpha^*+\alpha}{4}\biggr)
T_l\biggl(\frac{\alpha'+\alpha''}{4}\biggr) C_{-\delta\epsilon}\biggl(\frac{\phi_1-\phi_2}{2}\biggr) \biggr]=0.
\end{gather*}
From the last two equations, we conclude that
\[
T_l\biggl(\frac{\alpha^*-\alpha}{4}\biggr) S_{-\delta\epsilon}\biggl(\frac{\phi_2-\phi_1}{2}\biggr)-\delta\tau
T_l\biggl(\frac{\alpha'-\alpha''}{4}\biggr) S_{-\delta\epsilon}\biggl(\frac{\phi_2+\phi_1}{2}\biggr)=0,
\]
which is exactly the superposition formula given by \eqref{eq:superpos}.

It remains to prove that $\alpha^*$ defined by \eqref{eq:superpos} satisfies \eqref{eqorig1_4}. We are assuming that
$\alpha'$ and $\alpha''$ satisfy \eqref{eq:btanalytic_v1} and \eqref{eq:btanalytic_v2}, respectively and $\alpha$ satisfies \eqref{eqorig1_4}. Theorem \ref{thBT} implies that $\alpha'$ and $\alpha''$ also satisfy \eqref{eqorig1_4}. We introduce the following notation:
\begin{gather}\label{Sigmaetall}
Z:= \frac{S_{-\delta\epsilon}\bigl(\frac{\phi_2+\phi_1}{2}\bigr) } {S_{-\delta\epsilon}\bigl(\frac{\phi_2-\phi_1}{2}\bigr)},\qquad D:=\alpha''-\alpha', \qquad \Sigma:=\alpha''+\alpha',
\qquad G:= -\delta\tau Z T_l \biggl(\frac{D}{4}\biggr).
\end{gather}
It follows from \eqref{eq:superpos} that
\begin{equation*}
\alpha^*_{x_ix_i}=\alpha_{x_ix_i}+\frac{4}{1+l G^2}\biggl(\frac{-2l G}{1+l G^2}(G_{x_i})^2+G_{x_ix_i}\biggr), \qquad i=1,2.
\end{equation*}
Therefore, proving that $\alpha^*$ satisfies \eqref{eqorig1_4} is equivalent to
proving
\begin{equation}\label{eqalst}
\frac{4}{1+l G^2}
\biggl(\frac{-2l G}{1+l G^2}\bigl[(G_{x_1})^2-(G_{x_2})^2\bigr]+G_{x_1x_1}-G_{x_2x_2} \biggr) +\epsilon\delta ( S_l(\alpha^*)-S_l(\alpha))=0,
\end{equation}
where we used the fact that $\alpha$ is a solution of \eqref{eqorig1_4}.

Observe that
\[
S_l(\alpha^*)-S_l(\alpha)=2S_l\biggl(\frac{\alpha^*-\alpha}{2}\biggr) C_l\biggl(\frac{\alpha^*+\alpha}{2}\biggr)=
\frac{4G}{1+l G^2}C_l\biggl(\frac{\alpha^*+\alpha}{2}\biggr),
\]
and it follows from \eqref{eq:superpos} and a long computation that
\[
C_l\biggl(\frac{\alpha^*+\alpha}{2}\biggr)=\frac{1}{1+l G^2}\bigl[ \bigl(1-l G^2\bigr)C_l(\alpha)-2l G S_l(\alpha)\bigr].
\]
Therefore, \eqref{eqalst} reduces to
\begin{gather}
-2l\bigl[(G_{x_1})^2-(G_{x_2})^2\bigr]+\frac{\bigl(1+l G^2\bigr)}{G}(G_{x_1x_1}-G_{x_2x_2}) \nonumber\\
\qquad{}+\epsilon\delta\bigl[ \bigl(1-l G^2\bigr)C_l(\alpha)-2l G S_l(\alpha)\bigr]=0.\label{eqalst1}
\end{gather}


Using the expression of $G$, we have
\[
G_{x_i}=\frac{G D_{x_i}}{2S_l\bigl(\frac{D}{2}\bigr)},
\quad
G_{x_ix_i}=\frac{G}{2S_l\bigl(\frac{D}{2}\bigr)}\biggl( D_{x_ix_i}+\frac{l}{2}T_l\biggl(\frac{D}{4}\biggr) (D_{x_i})^2 \biggr), \qquad i=1,2.
\]
Moreover,
\[
\frac{l}{4S_l\bigl(\frac{D}{2}\bigr)}\biggl(\frac{-2G^2}{S_l\bigl(\frac{D}{2}\bigr)}+\bigl(1+l G^2\bigr)
T_l\biggl(\frac{D}{4}\biggr)\biggr)= \frac{-l\bigl(Z^2-1\bigr)}{8\bigl(C_l\bigl(\frac{D}{4}\bigr)\bigr)^2}.
\]
Substituting these expressions into \eqref{eqalst1}, we conclude that we need to prove that
\begin{gather}
\frac{-l\bigl(Z^2-1\bigr)}{8C_l^2\bigl(\frac{D}{4}\bigr)} \bigl[(D_{x_1})^2-(D_{x_2})^2\bigr]+
\frac{1+l G^2}{2S_l\bigl(\frac{D}{2}\bigr)}(D_{x_1x_1}-D_{x_2x_2})\nonumber\\
\qquad{}+\epsilon\delta \bigl[ \bigl(1-l G^2\bigr)C_l(\alpha)-2l G S_l(\alpha)\bigr]=0.\label{eqalst3}
\end{gather}

Since $\alpha'$ and $\alpha''$ satisfy \eqref{eqorig1_4}, it follows that
\begin{equation}\label{eqdifDxixi}
D_{x_1x_1}-D_{x_2x_2}=-2\epsilon\delta S_l\biggl(\frac{D}{2}\biggr) C_l\biggl(\frac{\Sigma}{2}\biggr),
\end{equation}
where $\Sigma$ is given by \eqref{Sigmaetall}. Moreover, since $\alpha'$ and $\alpha''$ satisfy
\eqref{eq:btanalytic_v1} and \eqref{eq:btanalytic_v2} respectively,
 a long but straightforward computation shows that
\begin{gather*}
(D_{x_1})^2-(D_{x_2})^2=\frac{-4}{\lambda_1^2\lambda_2^2}
\biggl[ -\biggl( C_l^2\Bigl(\frac{\alpha}{2}\Bigr)S_l^2\biggl(\frac{\alpha'}{2}\biggr)-
C_l^2\biggl(\frac{\alpha'}{2}\biggr)S_l^2\Bigl(\frac{\alpha}{2}\Bigr)\biggr) (1-\Lambda_1^2)\lambda_2^2 \\
\qquad\qquad{}-\biggl( C_l^2\Bigl(\frac{\alpha}{2}\Bigr)S_l^2\biggl(\frac{\alpha''}{2}\biggr)-
C_l^2\biggl(\frac{\alpha''}{2}\biggr)S_l^2\Bigl(\frac{\alpha}{2}\Bigr)\biggr) (1-\Lambda_2^2)\lambda_1^2 \\
\qquad\qquad{}+2\biggl( C_l^2\Bigl(\frac{\alpha}{2}\Bigr)S_l\biggl(\frac{\alpha'}{2}\biggr)S_l\biggl(\frac{\alpha''}{2}\biggr) - C_l\biggl(\frac{\alpha'}{2}\biggr)C_l\biggl(\frac{\alpha''}{2}\biggr)S_l^2\Bigl(\frac{\alpha}{2}\Bigr) \biggr)(1-\Lambda_1\Lambda_2)\lambda_1\lambda_2 \\
\qquad\qquad{} -S_l(\alpha)S_l\biggl(\frac{D}{2}\biggr) (\Lambda_2-\Lambda_1)\delta\tau\lambda_1\lambda_2 \biggr].
\end{gather*}
Therefore, it follows from the fact
\[
1-\Lambda_i^2=-\epsilon\delta\lambda_i^2, \qquad
1-\Lambda_1\Lambda_2=\frac{-\epsilon\delta\bigl(Z^2+1\bigr)}{Z^2-1} \lambda_1\lambda_2, \qquad
\Lambda_2-\Lambda_1=\frac{2\epsilon\delta Z}{Z^2-1}\lambda_1\lambda_2,
\]
 that
\begin{gather*}
(D_{x_1})^2-(D_{x_2})^2=\frac{-16\epsilon\delta}{Z^2-1}\biggl[ \biggl( C_l^2\Bigl(\frac{\alpha}{2} \Bigr) C_l^2\biggl(\frac{\Sigma}{4}\biggr)- S_l^2\Bigl(\frac{\alpha}{2}\Bigr) S_l^2\biggl(\frac{\Sigma}{4}\biggr) \biggr)
S_l^2\biggl(\frac{D}{4}\biggr)Z^2 \\
\hphantom{(D_{x_1})^2-(D_{x_2})^2=}{}
-\biggl(\! C_l^2\Bigl(\frac{\alpha}{2}\Bigr)S_l^2\biggl(\frac{\Sigma}{4}\biggr)-
S_l^2\Bigl(\frac{\alpha}{2} \Bigr) C_l^2\biggl(\frac{\Sigma}{4}\biggr)\!\biggr)
C_l^2\biggl(\!\frac{D}{4}\!\biggr)
 -\frac{1}{2}S_l(\alpha)S_l\biggl(\!\frac{D}{2}\!\biggr)\delta\tau Z
 \biggr],
\end{gather*}
where we used the fact that $S_l^2\bigl(\frac{D}{4}\bigr)Z^2=G^2 C_l^2\bigl(\frac{D}{4}\bigr)$.

Substituting this expression and \eqref{eqdifDxixi} into \eqref{eqalst3}, we obtain
a second degree polynomial in~$G$, whose coefficients vanish identically.
\end{proof}

\end{document}
